%==============================================================================
%  MATH 231 -- Linear Algebra I, Section 10 (Queens College, CUNY -- Fall 2026)
%  COURSE SYLLABUS
%  Instructor: Kevin Bedoya
%
%  FILENAME is the department's required format: number_section_lastname, i.e.
%  231_sec10_bedoya.pdf.  Do not rename the exported PDF.
%
%  SOURCES this document was assembled from, so that it can be re-checked:
%    - welcome-letter/welcome-letter.pdf   (voice, logistics, grade weights,
%      platforms, textbook stance, programming stance)
%    - schedule/MATH231_Fall2026_schedule.pdf (28 meetings, the two lost days,
%      the final-exam window)
%    - tentative_textbooks/                (the four texts and their free links)
%    - lectures/primer/                    (the pre-course foundations review)
%    - "Instructor Guide.pdf" (Math Dept)  (VOE, exam/homework guidance, and the
%      full contact list reproduced in Sec. 13)
%    - https://www.qc.cuny.edu/cetll/whats-in-a-syllabus/
%    - https://www.qc.cuny.edu/provost/effective-syllabus/
%    - https://www.qc.cuny.edu/provost/course-processes/
%    - QC undergraduate course listing (catalog description, hours, prereq)
%
%  THREE THINGS TO SET BEFORE HANDING THIS OUT.  Each is a single macro below:
%    \officehours  -- currently a placeholder; announce at the first meeting.
%    \midtermone / \midtermtwo -- tentative placements, see Sec. 7.
%  Everything else is sourced.
%
%  OSDA, NOT SPSV.  The office formerly called Special Services (SPSV) is now
%  the Office for Students with Disabilities and Access (OSDA), and its address
%  is OSDA@qc.cuny.edu.  The CETLL syllabus page still shows the old
%  qc.spsv@qc.cuny.edu; do not copy it back in.
%
%  ATTENDANCE.  Two different things are deliberately kept apart in Sec. 6: VOE
%  attendance for the first three weeks, which is a registration matter, and
%  grading, which does not use attendance -- per the Provost's guidance that
%  attendance may not be used to evaluate students.
%
%  Compile with pdflatex.  Standard TeX Live packages only.
%==============================================================================
\documentclass[11pt]{article}

\usepackage[T1]{fontenc}
\usepackage[utf8]{inputenc}
\usepackage{lmodern}
\usepackage{amsmath,amssymb}
\usepackage[margin=1in]{geometry}
\usepackage{microtype}
\usepackage{enumitem}
\usepackage{xcolor}
\usepackage{booktabs}
\usepackage{array}
\usepackage{longtable}
\usepackage{titlesec}
\usepackage{needspace}
\usepackage[hidelinks]{hyperref}
\usepackage{xurl}

\definecolor{accent}{HTML}{1F4E79}
\definecolor{rule}{HTML}{C8D3DF}
\definecolor{softbg}{HTML}{F2F6FA}
\definecolor{warnbg}{HTML}{FDF3E7}
\definecolor{warnrule}{HTML}{D98B3A}

\hypersetup{
  colorlinks=true,
  linkcolor=accent,
  urlcolor=accent,
  pdftitle={MATH 231 Linear Algebra I, Section 10 --- Syllabus, Fall 2026},
  pdfauthor={Kevin Bedoya}
}

\titleformat{\section}{\large\bfseries\color{accent}}{\thesection.}{0.5em}{}
\titlespacing*{\section}{0pt}{1.5em}{0.5em}
\titleformat{\subsection}{\normalsize\bfseries\color{accent}}{\thesubsection.}{0.5em}{}
\titlespacing*{\subsection}{0pt}{1.0em}{0.35em}

\setlist[itemize]{leftmargin=1.4em, itemsep=0.25em, topsep=0.35em}
\setlist[enumerate]{leftmargin=1.6em, itemsep=0.25em, topsep=0.35em}
\setlength{\parskip}{0.55em}
\setlength{\parindent}{0pt}

% Shaded call-out boxes.  The content is collected into a savebox first because
% \colorbox needs its whole argument at once and cannot straddle the
% \begin/\end split of an environment.
\newsavebox{\qcbox}
\newenvironment{callout}[1]
  {\par\medskip\needspace{5\baselineskip}%
   \begin{lrbox}{\qcbox}%
   \begin{minipage}{\dimexpr\textwidth-2\fboxsep-4pt}%
   \setlength{\parskip}{0.4em}%
   \textbf{\color{accent}#1}\par}
  {\end{minipage}\end{lrbox}%
   \begin{center}\colorbox{softbg}{\usebox{\qcbox}}\end{center}\medskip}

\newenvironment{alertbox}[1]
  {\par\medskip\needspace{5\baselineskip}%
   \begin{lrbox}{\qcbox}%
   \begin{minipage}{\dimexpr\textwidth-2\fboxsep-4pt}%
   \setlength{\parskip}{0.4em}%
   \textbf{\color{warnrule}#1}\par}
  {\end{minipage}\end{lrbox}%
   \begin{center}\colorbox{warnbg}{\usebox{\qcbox}}\end{center}\medskip}

% One MQR outcome: number, the official wording, the course-specific gloss,
% and where it is assessed.  A four-column table made these unreadably narrow.
\newcommand{\mqr}[4]{%
  \par\medskip\needspace{4\baselineskip}%
  {\bfseries\color{accent}MQR #1.}\ \emph{#2}\par\nobreak\vspace{0.2em}
  \begingroup\leftskip=1.5em\relax
  #3\par\nobreak\vspace{0.15em}
  {\footnotesize\itshape Assessed by: #4.}\par
  \endgroup}

%--------------------------------------------------------- fill these in -----
\newcommand{\officehours}{\emph{to be announced at the first class meeting and
  posted on the class Discord}}
\newcommand{\midtermone}{Thursday, October 8}
\newcommand{\midtermtwo}{Thursday, November 12}

\begin{document}

\begin{center}
  {\LARGE\bfseries\color{accent} MATH 231: Linear Algebra I}\\[0.4em]
  {\large\bfseries Section 10 \quad$\cdot$\quad Fall 2026}\\[0.5em]
  {\normalsize Department of Mathematics \quad$\cdot$\quad
   Queens College, City University of New York}
\end{center}

\vspace{0.3em}
{\color{rule}\hrule height 1.2pt}
\vspace{0.9em}

\begin{center}
\renewcommand{\arraystretch}{1.25}
\begin{tabular}{@{}>{\bfseries}l l@{}}
\toprule
Instructor      & Kevin Bedoya \\
Email           & \href{mailto:Kevin.Bedoya32@stu-mail.qc.cuny.edu}{Kevin.Bedoya32@stu-mail.qc.cuny.edu} \\
Meetings        & Tuesdays \& Thursdays, 12:10--2:00 PM \\
Location        & Kiely Hall 326 \\
Mode            & In-person \\
Office hours    & \officehours \\
Office location & Kiely Hall 331 (the Math Lab) \\
Credits         & 4 hours; 4 credits \\
Prerequisite    & MATH 141 or MATH 151, minimum grade C$-$ \\
Semester        & August 28 -- December 21, 2026 \\
\bottomrule
\end{tabular}
\end{center}

\begin{callout}{Email policy}
Write from your QC email account and put \texttt{MATH 231} in the subject line.
I answer within one business day on weekdays; replies over the weekend are
likely but not promised. Questions that other students would benefit from ---
about homework, a definition, an assignment deadline --- are better asked in
the class Discord, where everyone can see the answer.
\end{callout}

%==============================================================================
\section{Course description}
%==============================================================================

\textbf{Catalog description.} \emph{An introduction to linear algebra with
emphasis on techniques and applications. Topics to be covered include solutions
of systems of linear equations, vector spaces, bases and dimension, linear
transformations, matrix algebra, determinants, eigenvalues, and inner products.
Not open to students who are enrolled in or who have completed MATH 237.
Students who fail or withdraw from this course multiple times may be prohibited
from majoring in the sciences or mathematics; see the bulletin language for
your major.}

\textbf{How this section is taught.} Most students in this room are computer
science majors or in applied-mathematics-adjacent programs, and the section is
built accordingly. We cover the full catalog content, and we lean into the
computational side of it: how much a calculation actually costs, how a machine
really stores a number, and how the standard algorithms of data science are
built out of the vector and matrix operations we develop. Along the way we will
meet $k$-means clustering, $k$-nearest neighbours, the Gram--Schmidt process,
gradient descent, Google's PageRank, and image compression.

\textbf{On programming.} Our language is Python, used for homework, projects
and extra credit. \textbf{Your programming ability will not be assessed on
exams.} If you have never written a line of code, assignments come with
walkthroughs covering the syntax and machinery you need; if you have taken
CSCI 111, 211 or 212, you will be comfortable quickly. Projects are
implementation-based: you will be asked to implement an algorithm that has
already been developed in lecture. We will use \texttt{numpy}, \texttt{scipy}
and \texttt{matplotlib}.

%==============================================================================
\section{Pathways learning outcomes}
%==============================================================================

MATH 231 satisfies the CUNY Pathways \textbf{Mathematical and Quantitative
Reasoning} (MQR) requirement, and must meet all six MQR outcomes. Each is
stated below in its official wording, followed by what it looks like concretely
in this course and where it is assessed.

\mqr{1}{Interpret and draw appropriate inferences from quantitative
  representations, such as formulas, graphs, or tables.}
 {Read a matrix as a table of coefficients and recover the linear system it
  encodes; read an operation-count table and a growth plot to decide which of
  two algorithms scales better; read a table of distances to classify a point;
  read the spacing diagram of a floating-point system to say what precision is
  available near a given magnitude.}
 {homework, exams}

\mqr{2}{Use algebraic, numerical, graphical, or statistical methods to draw
  accurate conclusions and solve mathematical problems.}
 {Solve linear systems by elimination; compute dot products, norms,
  projections, determinants and eigenvalues; run the Gram--Schmidt process to
  produce an orthonormal basis; execute Lloyd's algorithm and the
  $k$-nearest-neighbours rule both by hand and in code.}
 {homework, projects, exams}

\mqr{3}{Represent quantitative problems expressed in natural language in a
  suitable mathematical format.}
 {Turn ``find the coordinates of this vector against this basis'' into the
  matrix equation $V\mathbf{x}=\mathbf{u}$; turn a described collection of
  objects into feature vectors in $\mathbb{R}^{n}$; turn a stated set of
  simultaneous constraints into a linear system and its matrix.}
 {homework, projects, exams}

\mqr{4}{Effectively communicate quantitative analysis or solutions to
  mathematical problems in written or oral form.}
 {Written homework solutions that justify each step rather than only stating an
  answer; a written report accompanying each project explaining the method, the
  results, and the limits of both; explaining your reasoning aloud in class and
  in office hours.}
 {homework, projects}

\mqr{5}{Evaluate solutions to problems for reasonableness using a variety of
  means, including informed estimation.}
 {Substitute a computed solution back into the original system to check it; use
  a Big-O estimate to judge in advance whether a computation is feasible at a
  given input size; use relative error and machine precision to decide how many
  digits of a numerical answer can be trusted; sanity-check a clustering
  against what the data plainly looks like.}
 {homework, projects, exams}

\mqr{6}{Apply mathematical methods to problems in other fields of study.}
 {Apply linear algebra to problems in computer science and data science:
  nearest-neighbour search and recommendation, clustering, image compression,
  PageRank, and the least-squares fitting that underlies regression.}
 {projects}

%==============================================================================
\section{Materials}
%==============================================================================

\textbf{No purchases are required.} Free online copies of every text are linked
below and collected in the course GitHub repository. Lectures do not follow any
single textbook --- I have written my own progression for this course. The
texts supply additional reading, alternate perspectives, and problems.

\renewcommand{\arraystretch}{1.3}
\begin{longtable}{@{} >{\raggedright\arraybackslash}p{6.4cm} >{\raggedright\arraybackslash}p{4.0cm} >{\raggedright\arraybackslash}p{5.0cm} @{}}
\toprule
\textbf{Text} & \textbf{Role in the course} & \textbf{Free access} \\
\midrule
\endhead
\bottomrule
\endlastfoot
Boyd \& Vandenberghe, \emph{Introduction to Applied Linear Algebra: Vectors,
Matrices, and Least Squares} (Cambridge, 2018)
 & Applied spine
 & \url{https://web.stanford.edu/~boyd/vmls/vmls.pdf} \\
Trefethen \& Bau, \emph{Numerical Linear Algebra} (SIAM, 1997)
 & Numerical methods, conditioning, QR, SVD
 & \url{https://www.stat.uchicago.edu/~lekheng/courses/309/books/Trefethen-Bau.pdf} \\
Strang, \emph{Linear Algebra and Its Applications}, 4th ed.\ (Cengage, 2006)
 & Subspaces and geometric intuition
 & Linked in the course repository \\
Axler, \emph{Linear Algebra Done Right}, 4th ed.\ (Springer, 2024)
 & Theoretical foundation
 & \url{https://linear.axler.net} \\
\end{longtable}

\subsection{Foundations primer --- read this first}

A self-study primer, \emph{MATH 231 Foundations Review}, is posted in the course
repository under \texttt{lectures/primer/}. It contains no new linear algebra:
it collects the algebra, notation, summation and calculus you already have, so
that we can speak a common language on day one. \textbf{Read it before or during
the first week.} If a page feels like review, skim it; if a page feels new, slow
down and work the examples with a pencil.

\subsection{Platforms and setup}

Everything below is free.

\renewcommand{\arraystretch}{1.3}
\begin{center}
\begin{longtable}{@{} >{\bfseries}l >{\raggedright\arraybackslash}p{11.4cm} @{}}
\toprule
\endhead
\bottomrule
\endlastfoot
Discord & Primary communication. \textbf{All announcements go through the class
Discord} --- schedule changes, homework, extra credit, projects, exams, and
general information. Invite link distributed by email before the first meeting. \\
GitHub & Lecture notes, homework, projects, practice exams, textbook links: \newline
\url{https://github.com/KevinB88/MATH231-Queens-College-Fall-2026-Linear-Algebra-I-} \\
Gradesly & Where grades are posted: \url{https://gradesly.com/}. You will sign
up with your CUNY ID and a randomly generated passcode; setup instructions are
distributed in week 1. \\
Git + GitHub account & How you submit work. You will push homework, projects and
extra credit to your own MATH 231 repository. A setup walkthrough --- creating an
account, installing git, structuring your directory, and what git actually is ---
is provided in week 1. \\
Python + an IDE & VS Code or PyCharm. Both are free. \\
\end{longtable}
\end{center}

\begin{alertbox}{We are not using Brightspace}
Course materials, announcements and grades live on GitHub, Discord and Gradesly
respectively. Do not wait for a Brightspace post; you will miss things. If you
are not in the class Discord within the first week, email me.
\end{alertbox}

%==============================================================================
\section{Course progression (tentative)}
%==============================================================================

The order below is the plan and is \textbf{tentative}. Pacing will be adjusted
to the room, and the authoritative version of what has been covered is always
the lecture notes posted in the repository. Announcements of any change go
through Discord.

\renewcommand{\arraystretch}{1.3}
\begin{center}
\begin{longtable}{@{} >{\bfseries}c >{\raggedright\arraybackslash}p{4.5cm} >{\raggedright\arraybackslash}p{9.2cm} @{}}
\toprule
\textbf{Unit} & \textbf{Subject} & \textbf{Contents} \\
\midrule
\endfirsthead
\toprule
\textbf{Unit} & \textbf{Subject} & \textbf{Contents} \\
\midrule
\endhead
\bottomrule
\endlastfoot
0 & Foundations primer & Self-study. Notation, algebra, summation, functions,
    the calculus assumed from MATH 141. \\
1 & Vectors & Vectors in $\mathbb{R}^{2}$ and $\mathbb{R}^{n}$; addition and
    scaling; the dot product, angle and projection; norms, metrics and
    Cauchy--Schwarz; span, basis, coordinates and dimension; linear
    independence; orthogonality and subspaces. \\
2 & Computational foundations & Counting operations; asymptotic analysis and
    Big-O; floating-point arithmetic, machine precision and error;
    $k$-nearest neighbours; $k$-means clustering; the Gram--Schmidt process;
    matrix multiplication as an algorithm. \\
3 & Matrices and linear systems & Linear systems from the coordinate problem;
    the matrix and the matrix--vector product; Gaussian elimination; matrix
    algebra, the transpose and the inverse; the determinant and the trace. \\
4 & Linear maps and subspaces & Linear transformations and their matrices;
    kernel and image; rank--nullity; the four fundamental subspaces;
    change of basis. \\
5 & Orthogonality and least squares & Orthogonal projection onto a subspace;
    the normal equations; least-squares fitting; the QR factorization. \\
6 & Eigenvalues and applications & Eigenvalues and eigenvectors;
    diagonalization; an introduction to the singular value decomposition;
    applications including PageRank and image compression. \\
\end{longtable}
\end{center}

Complex numbers and quaternions appear as enrichment material alongside Unit 1
and are not examinable unless announced otherwise.

%==============================================================================
\section{Grading}
%==============================================================================

\begin{center}
\renewcommand{\arraystretch}{1.3}
\begin{tabular}{@{} l c >{\raggedright\arraybackslash}p{7.6cm} @{}}
\toprule
\textbf{Component} & \textbf{Weight} & \textbf{Notes} \\
\midrule
Homework (5 assignments) & 20\% & Submitted via your GitHub repository. \\
Projects (3)             & 15\% & Python implementations with a written report. \\
Midterm 1                & 20\% & In class, in person. \\
Midterm 2                & 20\% & In class, in person. \\
Final examination        & 25\% & Cumulative; during the December 15--21 period. \\
\midrule
\textbf{Total}           & \textbf{100\%} & \\
\bottomrule
\end{tabular}
\end{center}

Extra credit will be offered during the semester and announced on Discord.

\subsection{Grade scale}

\begin{center}
\renewcommand{\arraystretch}{1.2}
\begin{tabular}{@{} l l l l l l @{}}
\toprule
A+ & 97--100 & B+ & 87--89.9 & C+ & 77--79.9 \\
A  & 93--96.9 & B  & 83--86.9 & C  & 73--76.9 \\
A$-$ & 90--92.9 & B$-$ & 80--82.9 & C$-$ & 70--72.9 \\
D+ & 67--69.9 & D & 63--66.9 & D$-$ & 60--62.9 \\
F  & below 60 & & & & \\
\bottomrule
\end{tabular}
\end{center}

\subsection{Late work and make-ups}

\begin{itemize}
  \item \textbf{Homework and projects} are accepted up to 48 hours late with a
        10\% penalty per day. After 48 hours the grade is zero unless you have
        arranged something with me \emph{in advance}.
  \item \textbf{Exams} are in class and cannot be made up except for a
        documented emergency, a religious observance (\S9.7), or an approved
        accommodation (\S9.3). Contact me as far in advance as you can; for an
        emergency, as soon as you are able.
  \item Talk to me early. A problem raised before a deadline has many more
        solutions than the same problem raised after it.
\end{itemize}

%==============================================================================
\section{Attendance and Verification of Enrollment}
%==============================================================================

\begin{alertbox}{Verification of Enrollment (VOE) --- the first three weeks}
CUNY requires me to take attendance at \textbf{every meeting during the first
three weeks of the semester} and to submit Verification of Enrollment in
CUNYfirst. For this section that means meetings 1 through 6:
\textbf{September 1, 3, 8, 10, 15 and 17}.

If you do not attend during this window, you may be reported as not actively
participating and \textbf{dropped from the course}. This is also how CUNY
verifies enrollment for \textbf{financial aid}, so a missed verification can
affect your aid disbursement. If something prevents you from attending during
the first three weeks, email me immediately --- before the fact if at all
possible.
\end{alertbox}

\textbf{After the VOE window.} Attendance is not itself a graded component of
this course. What is graded is listed in \S5. That said, this course moves
quickly and builds on itself; the lecture notes are thorough but they are not a
substitute for being in the room, and consistent absence reliably shows up in
exam scores. If you must miss a class, get the notes, read them, and come to
office hours with questions.

\textbf{Lateness.} Class starts at 12:10. Please arrive on time; if you do
arrive late, come in quietly and take a seat rather than waiting outside.

%==============================================================================
\section{Examinations}
%==============================================================================

There are three examinations: two midterms and a cumulative final. All are
taken \textbf{in class, in person}. There are no take-home exams. Programming
is not assessed on exams (\S1).

\begin{center}
\renewcommand{\arraystretch}{1.3}
\begin{tabular}{@{} l l >{\raggedright\arraybackslash}p{7.4cm} @{}}
\toprule
\textbf{Exam} & \textbf{Date} & \textbf{Coverage} \\
\midrule
Midterm 1 & \midtermone\ \emph{(tentative)}
  & Units 1--2: vectors, inner products, norms, basis and span, and the
    computational material through Gram--Schmidt. \\
Midterm 2 & \midtermtwo\ \emph{(tentative)}
  & Units 3--4: linear systems, matrix algebra, determinants, linear maps and
    subspaces. \\
Final     & During December 15--21
  & Cumulative, with emphasis on Units 5--6. \\
\bottomrule
\end{tabular}
\end{center}

\begin{callout}{How exam dates are confirmed}
The two midterm dates above are \textbf{tentative}. Each will be confirmed on
Discord and in class \textbf{at least two weeks in advance}, and once confirmed
it will not move. The final examination is scheduled by the College within the
December 15--21 period; our meeting slots in that window are Tuesday,
December 15 and Thursday, December 17. The official date and time will be
announced as soon as the College publishes the final examination schedule.
\end{callout}

Exams are designed so that a student who attends regularly, does the homework,
and studies can do well. Questions will not be recycled verbatim from homework
or review sheets: the goal is to test whether you understand the concepts, not
whether you have memorized specific worked examples.

%==============================================================================
\section{Homework, projects and submission}
%==============================================================================

\textbf{Five homework assignments} and \textbf{three projects} across the
semester. Both are submitted by committing and pushing to your own MATH 231
GitHub repository; a walkthrough covering git, GitHub accounts and directory
structure is provided in week 1.

\begin{itemize}
  \item \textbf{Homework} is a mix of hand computation and short proofs, with
        some Python. Show your reasoning --- an unjustified correct answer
        earns partial credit at best. This is MQR 4, and it is graded.
  \item \textbf{Projects} ask you to implement, in Python, an algorithm already
        developed in lecture, and to submit a short written report: what you
        implemented, what you observed, and what the limits of the method are.
  \item \textbf{Collaboration} is encouraged on homework --- talk to each other,
        form study groups, work problems together at the board. What you submit
        must be written by you, in your own words and your own code. Copying
        another student's solution, or letting yours be copied, is a violation
        of \S9.1.
\end{itemize}

%==============================================================================
\section{Course policies}
%==============================================================================

\subsection{Academic integrity}

\emph{Academic dishonesty is prohibited in the City University of New York and
is punishable by penalties, including failing grades, suspension, and
expulsion.} This course follows the CUNY Policy on Academic Integrity as adopted
by the CUNY Board of Trustees; the College's policy page is at
\url{https://www.qc.cuny.edu/cetll/academic-integrity}.

Violations include plagiarism, copying another student's work, unauthorized
collaboration on an exam, using prohibited materials during an exam, and
submitting work you cannot explain as your own. Suspected violations are
reported to the Office of Judicial Affairs. Queens College's Coordinator of
Judicial Affairs is Emanuel Avila
(\href{mailto:emanuel.avila@qc.cuny.edu}{emanuel.avila@qc.cuny.edu} or
\href{mailto:studentconduct@qc.cuny.edu}{studentconduct@qc.cuny.edu},
718-997-3971).

\subsection{Generative artificial intelligence}

You may use generative AI tools --- ChatGPT, Claude, GitHub Copilot and the like
--- on \textbf{homework and projects}, to learn Python syntax, to debug, to get
an explanation of a concept, or to check your understanding. I use these tools
myself and I would rather teach you to use them well than pretend they do not
exist. Two conditions:

\begin{enumerate}
  \item \textbf{Disclose it.} Say so in a comment in your code, or in your
        project README: which tool, and what you used it for. A sentence is
        enough.
  \item \textbf{Own it.} You must be able to explain every line you submit. I
        may ask you to.
\end{enumerate}

\textbf{Generative AI is not permitted on examinations.} Submitting AI-generated
work that you cannot explain, or failing to disclose AI use, is a violation of
\S9.1. Note also that these tools are confidently wrong about mathematics with
some regularity --- a habit of checking their output against what you know is
part of MQR 5.

\subsection{Students with disabilities and accommodations}

The \textbf{Office for Students with Disabilities and Access (OSDA)} supports
students with qualifying disabilities under the Americans with Disabilities Act
by arranging reasonable accommodations. If you have previously received
accommodations, believe you may have a disability, or have a temporary
disability, contact OSDA:

\begin{center}
\renewcommand{\arraystretch}{1.25}
\begin{tabular}{@{} >{\bfseries}l l @{}}
\toprule
Main office & Kiely Hall 104/108 \\
Phone       & 718-997-5870 \\
Email       & \href{mailto:OSDA@qc.cuny.edu}{OSDA@qc.cuny.edu} \\
Testing office & Kiely Hall 220, 718-997-3775,
                 \href{mailto:osda.testing@qc.cuny.edu}{osda.testing@qc.cuny.edu} \\
\bottomrule
\end{tabular}
\end{center}

\textbf{Accommodations are not retroactive}, so register sooner rather than
later. Once OSDA sends me your accommodation letter, come talk to me --- early
in the semester if you can, and well before an exam --- so we can make the
arrangements without a rush. You are never required to disclose the nature of a
disability to me.

\emph{Note: this office was formerly called the Office of Special Services
(SPSV). Older documents may still use that name or the old email address.}

\subsection{Student wellness}

College is hard, and it is harder when something else is going wrong. Anxiety,
difficulty concentrating, loss, illness and stressful life events all affect
academic performance, and none of them are things you have to handle alone.
Confidential mental-health services are available on campus through
\textbf{Counseling Services}, on the first floor of Frese Hall:
718-997-5420, \href{mailto:CounselingServices@qc.cuny.edu}{CounselingServices@qc.cuny.edu},
\url{https://www.qc.cuny.edu/cs/}.

\textbf{A 24/7 crisis counselor is available: text ``CUNY'' to 741741.}

If something is going on that is affecting your work in this class, you are
welcome to tell me as much or as little as you want to. I would much rather hear
about it early than reconstruct it from a transcript at the end of the semester.

\subsection{Course withdrawals and incomplete grades}

Deadlines for dropping and withdrawing are set by the College, not by me, and
the grades \texttt{W}, \texttt{WN}, \texttt{WU} and \texttt{WA} have different
consequences for your record and your financial aid. Check the deadlines
yourself and talk to me before you withdraw:
\url{https://www.qc.cuny.edu/provost/withdrawals-incomplete-grades/}.

An \texttt{INC} is possible only when most of the coursework is already complete
and something outside your control prevents you from finishing; it is arranged
in advance, never awarded automatically. CUNY's uniform grade glossary is at
\url{https://www.cuny.edu/about/administration/offices/ur/resources/registrars-staff/CUNY_UNIFORM_GRADE_GLOSSARY091109.pdf}.

\subsection{Repeating this course}

Be aware of the College's policy on repeating courses and F-grade replacement:
\url{https://www.qc.cuny.edu/aac/wp-content/uploads/sites/84/2022/04/Repeating-Courses-and-the-F-Grade-Replacement-Policy.pdf}.
As the catalog description notes, students who fail or withdraw from MATH 231
multiple times may be prohibited from majoring in the sciences or mathematics.

\subsection{Religious observance}

Queens College accommodates religious observance, including for examinations.
If a scheduled exam or deadline conflicts with a religious obligation, tell me
as far in advance as possible and we will arrange an alternative. The College's
policy is at \url{https://www.qc.cuny.edu/provost/religious-accommodations/}.

\subsection{Grade appeals}

If you believe a grade is in error, come to me first --- most of these are
arithmetic and take five minutes to fix. If we cannot resolve it, appeals go
through the Undergraduate Scholastic Standards Committee:
\url{https://www.qc.cuny.edu/saa/}.

\subsection{Use of student work}

Samples of student work are occasionally made available to accreditation
professionals for periodic program reviews. Anonymity is assured under these
circumstances.

\subsection{Course evaluations}

During the final four weeks of the semester you will be asked to complete an
online course evaluation. All responses are completely anonymous; identifying
information is erased once the evaluation has been submitted. I read these, and
they change how I teach --- please fill it out.

\subsection{Non-discrimination and privacy}

Queens College does not discriminate against students on the basis of
pregnancy, childbirth or related conditions. Your educational records are
protected under the Family Educational Rights and Privacy Act (FERPA); see
\url{https://www.qc.cuny.edu/provost/guide-to-ferpa/}.

\subsection{Emergency drills}

\textbf{A campus fire drill is scheduled for sometime during the week of
September 28 -- October 1, 2026.} Two of our meetings fall in that window
(Tuesday, September 29 and Thursday, October 1). If the alarm sounds during
class, we evacuate Kiely Hall together by the nearest stairwell --- not the
elevator --- and reassemble outside before returning. Any class time lost will
be made up.

%==============================================================================
\section{Getting help}
%==============================================================================

In rough order of how quickly you will get an answer:

\begin{enumerate}
  \item \textbf{The class Discord.} Ask in the open channel. Other students will
        often answer before I do, and answering someone else's question is one
        of the best ways to find out whether you actually understand something.
  \item \textbf{Office hours}, Kiely Hall 331. Come with a specific question or
        come with nothing and we will find one. You do not need permission and
        you are not interrupting.
  \item \textbf{Email}, per the policy on page 1.
  \item \textbf{The Math Lab}, Kiely Hall 331 --- free drop-in tutoring and a
        library of videos: \url{https://sites.google.com/view/qc-mathtutoring/home}.
  \item \textbf{The QC Learning Commons}, Kiely 131 and 144:
        \url{https://www.qc.cuny.edu/academics/qclc/}.
  \item \textbf{Form a study group.} Genuinely: the students who do well in
        courses like this one are disproportionately the ones who work
        together.
\end{enumerate}

%==============================================================================
\section{Important contacts}
%==============================================================================

\renewcommand{\arraystretch}{1.3}
\begin{longtable}{@{} >{\raggedright\arraybackslash}p{4.6cm} >{\raggedright\arraybackslash}p{10.9cm} @{}}
\toprule
\textbf{Office} & \textbf{Contact} \\
\midrule
\endfirsthead
\toprule
\textbf{Office} & \textbf{Contact} \\
\midrule
\endhead
\bottomrule
\endlastfoot

Mathematics Department
 & 718-997-5811 \newline
   General: \href{mailto:Math@qc.cuny.edu}{Math@qc.cuny.edu} \newline
   Registration: \href{mailto:MathRegistration@qc.cuny.edu}{MathRegistration@qc.cuny.edu} \newline
   Placement: \href{mailto:MathPlacement@qc.cuny.edu}{MathPlacement@qc.cuny.edu} \newline
   \emph{Write to these addresses only --- not to office staff individually.} \\

Office for Students with Disabilities and Access (OSDA)
 & Kiely Hall 104/108 \newline
   718-997-5870 \newline
   \href{mailto:OSDA@qc.cuny.edu}{OSDA@qc.cuny.edu} \\

OSDA Testing Office
 & Kiely Hall 220 \newline
   718-997-3775 \newline
   \href{mailto:osda.testing@qc.cuny.edu}{osda.testing@qc.cuny.edu} \\

Academic Integrity / Judicial Affairs
 & Emanuel Avila, Coordinator of Judicial Affairs \newline
   718-997-3971 \newline
   \href{mailto:emanuel.avila@qc.cuny.edu}{emanuel.avila@qc.cuny.edu} \ or \
   \href{mailto:studentconduct@qc.cuny.edu}{studentconduct@qc.cuny.edu} \\

Counseling Services
 & Frese Hall, 1st floor \newline
   718-997-5420 \newline
   \href{mailto:CounselingServices@qc.cuny.edu}{CounselingServices@qc.cuny.edu} \newline
   \url{https://www.qc.cuny.edu/cs/} \newline
   \textbf{24/7 crisis counselor: text ``CUNY'' to 741741} \\

Immigrant Student Support Services
 & 718-997-3990 \newline
   \href{mailto:ImmigrantSupport@qc.cuny.edu}{ImmigrantSupport@qc.cuny.edu} \newline
   \url{https://www.qc.cuny.edu/immi/} \\

Technology Services (ITS)
 & I-Building 200 \newline
   718-997-4444 \newline
   \href{mailto:support@qc.cuny.edu}{support@qc.cuny.edu} \newline
   Support tickets: \url{https://support.qc.cuny.edu/support/home} \\

Security Office
 & Main Gate \newline
   718-997-5911 / 718-997-5912 \\

Audio Visual (media equipment)
 & 718-997-5960 \\

Human Resources
 & Kiely Hall 163 \newline
   718-997-4455 \newline
   \href{mailto:Ohr.Employeeservices@qc.cuny.edu}{Ohr.Employeeservices@qc.cuny.edu} \\

Payroll Office
 & Kiely Hall 151 \newline
   718-997-5765 \newline
   \href{mailto:Ohr.Payroll@qc.cuny.edu}{Ohr.Payroll@qc.cuny.edu} \\

Provost --- course processes and policies
 & \url{https://www.qc.cuny.edu/provost/course-processes/} \\

CUNY academic calendar
 & \url{https://www.cuny.edu/academics/academic-calendars} \\

\end{longtable}

%==============================================================================
\section{Meeting schedule}
%==============================================================================

Twenty-eight regular lecture meetings, September 1 through December 10. Two
scheduled meetings are lost to the College calendar. \textbf{This schedule is
subject to change}; changes are announced on Discord.

\begin{center}
\renewcommand{\arraystretch}{1.15}
\begin{longtable}{@{} c l c @{\hspace{2em}} c l c @{}}
\toprule
\textbf{\#} & \textbf{Date} & \textbf{Day} &
\textbf{\#} & \textbf{Date} & \textbf{Day} \\
\midrule
\endhead
\bottomrule
\endlastfoot
1  & Sep 1  & Tue & 15 & Oct 22 & Thu \\
2  & Sep 3  & Thu & 16 & Oct 27 & Tue \\
3  & Sep 8  & Tue & 17 & Oct 29 & Thu \\
4  & Sep 10 & Thu & 18 & Nov 3  & Tue \\
5  & Sep 15 & Tue & 19 & Nov 5  & Thu \\
6  & Sep 17 & Thu & 20 & Nov 10 & Tue \\
7  & Sep 22 & Tue & 21 & Nov 12 & Thu \\
8  & Sep 24 & Thu & 22 & Nov 17 & Tue \\
9  & Sep 29 & Tue & 23 & Nov 19 & Thu \\
10 & Oct 1  & Thu & 24 & Nov 24 & Tue \\
11 & Oct 6  & Tue & -- & \emph{Nov 26} & \emph{Thu} \\
12 & Oct 8  & Thu & 25 & Dec 1  & Tue \\
-- & \emph{Oct 13} & \emph{Tue} & 26 & Dec 3  & Thu \\
13 & Oct 15 & Thu & 27 & Dec 8  & Tue \\
14 & Oct 20 & Tue & 28 & Dec 10 & Thu \\
\end{longtable}
\end{center}

\begin{center}
\renewcommand{\arraystretch}{1.25}
\begin{tabular}{@{} l l @{}}
\toprule
\textbf{Date} & \textbf{What happens} \\
\midrule
Sep 1--17 (meetings 1--6) & VOE attendance window (\S6) \\
Sep 28 -- Oct 1           & Campus fire drill, exact day to be announced (\S9.12) \\
\midtermone               & Midterm 1 \emph{(tentative --- confirmed two weeks ahead)} \\
Tue, Oct 13               & \textbf{No class} --- College follows a Monday schedule \\
\midtermtwo               & Midterm 2 \emph{(tentative --- confirmed two weeks ahead)} \\
Thu, Nov 26               & \textbf{No class} --- College closed, Thanksgiving \\
Thu, Dec 10               & Last regular lecture \\
Dec 15--21                & Final examination period \\
\bottomrule
\end{tabular}
\end{center}

\vspace{1.2em}
{\color{rule}\hrule height 1.2pt}
\vspace{0.6em}

\begin{center}
\footnotesize\itshape
This syllabus is subject to change. Any change will be announced in class and on
the course Discord, and the current version is always the one in the course
GitHub repository.\\[0.3em]
MATH 231 Section 10 \textbullet\ Fall 2026 \textbullet\ Kevin Bedoya
\end{center}

\end{document}
